\documentclass[10pt,letterpaper]{article}
\usepackage[ficha]{../latex_common/estilo_monitoria}
\usepackage{multicol}

\fancyhead[L]{\footnotesize\sffamily\color{GrisTexto}\textbf{Ficha Resumen 04 --- Patrones de Vectores y Matrices Estáticos}}
\fancyhead[R]{\footnotesize\sffamily\color{Turquesa}\textbf{Monitoría Algoritmos 30h}}

\begin{document}
\thispagestyle{empty}

\begin{tcolorbox}[enhanced,colback=AzulMarino,colframe=AzulMarino,arc=4pt,top=3.5mm,bottom=3.5mm,left=4mm,right=4mm,boxrule=0pt]
  \tcbox[on line,colback=Turquesa,colframe=Turquesa,arc=2pt,boxsep=1pt,left=4pt,right=4pt]{\footnotesize\sffamily\bfseries\color{white} FICHA IMPRIMIBLE 04}\hspace{6pt}%
  {\large\bfseries\sffamily\color{white} Patrones Maestros de Vectores y Matrices Estáticos}\\[2pt]
  {\footnotesize\sffamily\color{AzulSuave} Memoria fija definida con \texttt{const MAX}, recorridos, búsqueda, Burbuja, diagonales y paso a funciones}
\end{tcolorbox}

\begin{center}
  \includegraphics[width=0.49\textwidth,height=4.8cm,keepaspectratio]{mem_03_vector_estatico_1d.png}\hfill
  \includegraphics[width=0.49\textwidth,height=4.8cm,keepaspectratio]{mem_04_matriz_estatica_2d.png}
\end{center}

\vspace{-6pt}
\begin{peligro}{Regla Inquebrantable de Arreglos Estáticos}
\scriptsize
Declara SIEMPRE capacidad máxima constante: \codigo{const int MAX = 100; int v[MAX];} (o \codigo{int m[MAX\_F][MAX\_C];}) y valida \codigo{1 <= n \&\& n <= MAX}. \textbf{PROHIBIDO:} \codigo{int v[n];} (VLA), \codigo{std::vector}, \codigo{new}/\codigo{delete} o punteros. Índices válidos: \codigo{0} hasta \codigo{n - 1} (\textbf{nunca} \codigo{v[n]}).
\end{peligro}

\vspace{-4pt}
\begin{multicols}{2}
\begin{lstlisting}[style=cppmonitoria, caption={Patrones 1D: Máximo, Búsqueda y Burbuja}]
// 1. Mayor con posicion (inicializar en 0)
int posMax = 0;
for (int i = 1; i < n; i++) {
    if (v[i] > v[posMax]) posMax = i;
}

// 2. Busqueda secuencial (-1 si no esta)
int pos = -1;
for (int i = 0; i < n && pos == -1; i++) {
    if (v[i] == buscado) pos = i;
}

// 3. Ordenamiento Burbuja ascendente
for (int p = 0; p < n - 1; p++) {
    for (int j = 0; j < n - 1 - p; j++) {
        if (v[j] > v[j + 1]) {
            int tmp = v[j];
            v[j] = v[j + 1];
            v[j + 1] = tmp;
        }
    }
}
\end{lstlisting}

\columnbreak

\begin{lstlisting}[style=cppmonitoria, caption={Patrones 2D y Paso a Funciones}]
const int MAX_F = 10, MAX_C = 10;

// Matriz: 2da dimension [MAX_C] OBLIGATORIA
void mostrarMat(const int m[][MAX_C], int f, int c) {
    for (int i = 0; i < f; i++) {      // Fila fija
        for (int j = 0; j < c; j++) {  // Columnas
            cout << m[i][j] << "\t";
        }
        cout << endl;
    }
}

// Suma de una columna fija j
int sumaCol(const int m[][MAX_C], int f, int col) {
    int s = 0;
    for (int i = 0; i < f; i++) s += m[i][col];
    return s;
}

// Diagonales (solo matriz cuadrada n x n)
// Principal: m[i][i]
// Secundaria: m[i][n - 1 - i] (i + j == n - 1)
\end{lstlisting}
\end{multicols}

\vspace{-6pt}
\begin{consejo}{Resumen de Firmas en Funciones (\texttt{main} $\rightarrow$ Función)}
\scriptsize
\textbf{Vector:} firma \codigo{void leer(int v[], int n)} o \codigo{void ver(const int v[], int n)} $\Rightarrow$ llamada \codigo{ver(vec, n);} (sin corchetes).\\
\textbf{Matriz:} firma \codigo{void leerM(int m[][MAX\_C], int f, int c)} $\Rightarrow$ llamada \codigo{leerM(mat, f, c);} (sin corchetes).
\end{consejo}

\end{document}
